\chapter{Deep Dive: Files at Scale --- Presigned URLs, Endpoints, and Why Bytes Don't Belong in Postgres}
\label{ch:dd-storage}

Chapter~9 introduced MinIO as ``S3 at home'' and explained the
two-endpoints bug. That was enough to make submissions work. This chapter
asks what the design costs once the files are many, large, and important
--- the moment a student's graded essay is the only copy anywhere. The
project's own code is the specimen: one \code{S3Config}, one
\code{S3StorageService}, one column called \code{file\_object\_key}.

Three questions organize the chapter:

\begin{enumerate}
  \item \textbf{Where should bytes live} --- in the database, on a
  disk, or in an object store --- and what exactly does each choice cost?
  \item \textbf{Who should move the bytes} --- the browser through the
  service, or the browser straight to storage with a signed permission
  slip?
  \item \textbf{What happens when the two stores disagree} --- a row that
  names an object that does not exist, or an object no row names?
\end{enumerate}

\section{Why object storage, and what it solves}

A submission has two halves: facts about it (who, when, which attempt,
what score) and the file itself. The facts are relational and tiny; the
file is opaque and potentially 20\,MB. Three homes are possible for the
file.

\textbf{In Postgres, as \code{bytea}.} One store, one transaction, one
backup. It is also where the trouble starts. Every row read pulls the
blob through shared buffers; \code{pg\_dump} grows by the size of every
essay ever submitted; each update rewrites the whole value and leaves a
dead tuple for autovacuum; the write-ahead log carries every byte twice
(WAL, then the page); replication streams it to every standby; and a
download must be materialized in the JVM before it reaches the socket.
Postgres will do it, and up to a few gigabytes nobody notices. Past that,
every operational task --- backup, restore, replication lag, vacuum ---
is dominated by data the database adds no value to.

\textbf{On a volume mounted into the service.} Fast and simple until the
service has two replicas on two nodes, which is the first thing
Kubernetes wants to do. A \code{ReadWriteOnce} claim (Chapter~19) binds
to one node; \code{ReadWriteMany} needs NFS or a storage driver, and the
service now owns file layout, permissions, cleanup, and its own backup.
The disk is also inside the failure domain of the pod.

\textbf{In an object store.} Flat namespace, HTTP API, content-addressed
durability, and --- decisively --- the ability to hand a client a URL and
step out of the data path. Backups become a bucket replication rule.
Serving becomes a CDN's job. The database keeps a 500-character key.

\begin{decision}[a key in Postgres, bytes in MinIO]
The project stores metadata in the \code{submission} table and the file
under a key in the \code{ts-submissions} bucket. The row references the
object by name; nothing references the row from the object. The
alternatives lose on the axis that matters for a school: submissions are
written once, read a few times, and kept for years --- the profile object
storage is built for and databases are worst at. The decision would
reverse only if files were tiny (a few KB), needed transactional
atomicity with the row, and were read on every row access --- an avatar
thumbnail, say. Even then the honest answer is usually ``still S3, plus
a cache''.
\end{decision}

\section{Reading the project's code}

\subsection{Two clients, two endpoints}

The gateway to everything is \code{S3Config}. Read it with the compose
environment in mind (\code{APP\_S3\_ENDPOINT=http://minio:9000},
\code{APP\_S3\_PUBLIC\_ENDPOINT=http://localhost:9000}):

\begin{lstlisting}[language=Java]
@Bean
public S3Client s3Client() {
    return S3Client.builder()
            .endpointOverride(URI.create(endpoint))            (*@\co{1}@*)
            .region(Region.of(region))
            .credentialsProvider(credentialsProvider())
            .serviceConfiguration(S3Configuration.builder()
                    .pathStyleAccessEnabled(true).build())     (*@\co{2}@*)
            .build();
}

@Bean
public S3Presigner s3Presigner() {
    return S3Presigner.builder()
            .endpointOverride(URI.create(publicEndpoint))      (*@\co{3}@*)
            .region(Region.of(region))
            .credentialsProvider(credentialsProvider())
            .serviceConfiguration(S3Configuration.builder()
                    .pathStyleAccessEnabled(true).build())
            .build();
}
\end{lstlisting}

\begin{description}
  \item[\co{1}] The client that moves bytes talks to the
  \emph{internal} name. Inside Compose or the cluster, \code{minio} is a
  DNS name on the private network (Chapter~18); from a browser it is
  nothing.
  \item[\co{2}] Path-style addressing: the bucket goes in the URL path,
  so the request is \code{/ts-submissions/key} on host \code{minio:9000}.
  AWS's default is virtual-host style, where the bucket is a DNS
  subdomain --- host \code{ts-submissions.s3.amazonaws.com}, path
  \code{/key}. MinIO can do that too, but only
  when it knows its own domain (\code{MINIO\_DOMAIN}) and wildcard DNS
  exists for it; on a home network neither is true, so the path form is
  the only one that resolves.
  \item[\co{3}] The presigner never sends a request. It only computes
  signatures --- and it must compute them for the host the
  \emph{browser} will contact. This is a separate bean for that reason
  alone. Section~\ref{sec:sigv4} explains why the host cannot be swapped
  afterwards.
\end{description}

\subsection{The key, the upload, the download}

\code{S3StorageService} is short enough to read whole; the important
lines are these:

\begin{lstlisting}[language=Java]
@PostConstruct
void ensureBucketExists() {
    try {
        s3Client.headBucket(HeadBucketRequest.builder()
                .bucket(bucket).build());
    } catch (NoSuchBucketException e) {                     (*@\co{1}@*)
        s3Client.createBucket(CreateBucketRequest.builder()
                .bucket(bucket).build());
    }
}

public String upload(MultipartFile file, Long assignmentId,
                     Long studentId) {
    String key = "submissions/%d/%d/%s-%s".formatted(
            assignmentId, studentId, UUID.randomUUID(), original); (*@\co{2}@*)
    s3Client.putObject(PutObjectRequest.builder()
                    .bucket(bucket).key(key)
                    .contentType(file.getContentType())
                    .contentLength(file.getSize()).build(),
            RequestBody.fromInputStream(file.getInputStream(),
                    file.getSize()));                          (*@\co{3}@*)
    return key;
}

public String presignedDownloadUrl(String key) {
    GetObjectPresignRequest presignRequest = GetObjectPresignRequest
            .builder()
            .signatureDuration(Duration.ofMinutes(presignedExpiryMinutes))
            .getObjectRequest(GetObjectRequest.builder()
                    .bucket(bucket).key(key).build())
            .build();
    return s3Presigner.presignGetObject(presignRequest).url().toString();
}
\end{lstlisting}

\begin{description}
  \item[\co{1}] Only \code{NoSuchBucketException} is caught. A
  connection refused --- MinIO not up yet --- is an
  \code{SdkClientException}, which escapes, fails bean creation, and
  takes the whole service down. The Kubernetes manifest for MinIO says
  so in its first comment: course-service will \code{CrashLoopBackOff}
  until MinIO answers. Compose hides this with
  \code{depends\_on: service\_healthy}.
  \item[\co{2}] The key is the whole data model of an object store. S3
  has no directories; the slashes are characters in a string. But
  \emph{prefixes} are what you can list by, so
  \code{submissions/\{assignmentId\}/\{studentId\}/} means ``every attempt
  by this student for this assignment'' is one \code{ListObjectsV2} call
  with a prefix. The UUID guarantees a fresh key per attempt, so
  nothing is ever overwritten and versioning is not needed for safety.
  The original filename is appended for human readability in the
  console; the database keeps it separately in \code{file\_name} because
  a key is not a good place for spaces, Unicode, or \code{..}.
  \item[\co{3}] The bytes flow \emph{through} this JVM:
  Spring's multipart support has already parsed them (spooled to a temp
  file above a small threshold), and now they are streamed to MinIO. The
  request thread is occupied for the whole transfer, twice.
\end{description}

And the seam with the relational half, in \code{SubmissionService}:

\begin{lstlisting}[language=Java]
@Transactional
public SubmissionResponse submit(...) {
    ...
    if (hasFile) {
        String key = s3StorageService.upload(file, assignmentId,
                student.getId());                              (*@\co{1}@*)
        submission.setFileObjectKey(key);
        submission.setFileName(file.getOriginalFilename());
    }
    submission = submissionRepository.save(submission);        (*@\co{2}@*)
    return toResponse(submission);
}
\end{lstlisting}

\begin{description}
  \item[\co{1}] An HTTP call to another system, inside a database
  transaction. It is not, and cannot be, part of that transaction.
  \item[\co{2}] If this insert fails --- the unique constraint on
  \code{(assignment\_id, student\_id, attempt\_number)} loses a race
  between two tabs, or Postgres is briefly unreachable --- the
  transaction rolls back and the object uploaded at \co{1} stays in the
  bucket, named by a key no row will ever record. Chapter~11 called
  this the dual-write problem when the second system was Kafka; it is
  the same problem with a different second system.
\end{description}

\section{Signature Version 4, and why the host is welded in}
\label{sec:sigv4}

A presigned URL is an ordinary S3 request with the authentication moved
from headers into the query string. The project's download link looks
like this (wrapped):

\begin{lstlisting}
http://localhost:9000/ts-submissions/submissions/7/3/2f9c...-essay.pdf
  ?X-Amz-Algorithm=AWS4-HMAC-SHA256
  &X-Amz-Credential=minioadmin/20260907/us-east-1/s3/aws4_request
  &X-Amz-Date=20260907T091500Z
  &X-Amz-Expires=900
  &X-Amz-SignedHeaders=host
  &X-Amz-Signature=3c1f...e9a0
\end{lstlisting}

The signature is computed in two stages. First a \emph{signing key} is
derived from the secret through a chain of HMACs over the date, region,
service name and the literal \code{aws4\_request}, so the secret itself
never touches a request. Then that key signs a \emph{string to sign}
that contains the SHA-256 of a \emph{canonical request}: the HTTP
method, the canonical URI (bucket and key, path-style), the sorted query
string, the signed headers, and the payload hash. Look at
\code{X-Amz-SignedHeaders=host}: the \code{Host} header is always
signed. Change the host --- \code{minio} to \code{localhost}, or
\code{localhost} to \code{127.0.0.1} --- and the canonical request
changes, the hash changes, the signature no longer matches, and MinIO
answers \code{403 SignatureDoesNotMatch}. There is no way to sign against
the internal name and rewrite the link for the outside world. Whoever
signs must know the exact origin the browser will use.

\begin{handson}[break a signature on purpose]
Log in, submit an assignment with a file, and copy the
\code{downloadUrl} the API returns. It works as issued:
\begin{lstlisting}[language=bash]
$ curl -s -o /dev/null -w "%{http_code}\n" "$URL"
200
\end{lstlisting}
Now change nothing but the host, to an address that reaches the very
same MinIO container:
\begin{lstlisting}[language=bash]
$ curl -s "${URL/localhost/127.0.0.1}" | head -c 160
<?xml version="1.0" encoding="UTF-8"?><Error>
<Code>SignatureDoesNotMatch</Code><Message>The request signature we
calculated does not match the signature you provided.
\end{lstlisting}
Same server, same object, same credentials --- and a different
\code{Host} header is enough. That is the whole reason
\code{public-endpoint} exists.
\end{handson}

\begin{pitfall}[download links work in Compose and fail in the cluster]
Symptom: on the cluster, clicking a submission's file opens
\code{http://localhost:9000/...} and the browser shows connection
refused. Cause: the course-service Deployment sets
\code{APP\_S3\_ENDPOINT} to \code{http://minio:9000} but never sets
\code{APP\_S3\_PUBLIC\_ENDPOINT}, so the config-server default
\code{http://localhost:9000} is signed into every link --- an address
that meant ``the developer's laptop'' and now means ``the user's
laptop''. Fix in two parts: expose MinIO's S3 port through an Ingress
rule on its own host (path-style signatures include the path, so the
Ingress must not add a prefix), and set the public endpoint to that
origin:
\begin{lstlisting}[language=yaml]
- { name: APP_S3_PUBLIC_ENDPOINT,
    value: "https://files.ts-server.tailbfe002.ts.net" }
\end{lstlisting}
Until a second hostname exists, the honest interim is to stream
downloads through course-service (\code{GetObject} into the response
body) so the browser never needs to reach MinIO at all --- slower, but
correct.
\end{pitfall}

\section{Who moves the bytes}

Today's upload path is: browser $\rightarrow$ Traefik $\rightarrow$
api-gateway $\rightarrow$ course-service $\rightarrow$ MinIO. Every hop
sees every byte. The service enforces \code{max-file-size: 20MB} and
\code{max-request-size: 25MB} in its multipart settings, so the worst
case is bounded; but each upload holds a Tomcat thread and a temp file
for its whole duration, the reactive gateway (Chapter~4) forwards the
body chunk by chunk but still spends its event loop on it, and any
ingress-level body limit or idle timeout applies in front of all of
them. Under a deadline --- 200 students uploading at 23:59 --- the
service's thread pool is the bottleneck for work it does not need to do.

The alternative is the pattern Chapter~9 mentioned for downloads,
applied to uploads: the service signs a \emph{PUT} for a key it chose,
the browser uploads directly to MinIO, and only then does the service
record the row. The service spends microseconds per file and never sees
the bytes. What it gives up: it cannot inspect the content before it
lands (virus scan, type sniffing) and it must learn, from somewhere, that
the upload actually completed. Section~\ref{sec:presigned-put} builds
this properly.

Beyond about 100\,MB a single PUT is the wrong tool on any network:
\emph{multipart upload} splits the object into parts of at least
5\,MB, uploads them in parallel, resumes after a failure by re-sending
one part, and assembles them server-side in a final call. The SDK's
\code{S3TransferManager} does this automatically above a threshold. The
operational tail: abandoned uploads leave invisible parts that are
billed and counted --- every bucket should carry a lifecycle rule that
aborts incomplete multipart uploads after a few days.

\section{Pros and cons, candidly}

\begin{center}
\begin{tabular}{@{}p{0.3\linewidth}p{0.62\linewidth}@{}}
\toprule
\textbf{Strength} & \textbf{Where the project pays for it} \\
\midrule
Database stays small; backup is metadata only &
  Two stores to back up, and only one of them transactionally \\
Standard API: MinIO today, S3 tomorrow with two env vars &
  Two endpoints to keep straight, and the cluster gets one wrong \\
Presigned downloads take the service out of the data path &
  Uploads still go through the service, thread-per-file \\
Keys give free ``directories'' by prefix &
  Key design is a one-way door; renaming means copying \\
Durability is storage's problem &
  On this hardware it is one container and one 5\,GB claim \\
Credentials are two strings &
  They are the MinIO \emph{root} credentials, in every course-service pod \\
\bottomrule
\end{tabular}
\end{center}

The last row deserves emphasis. \code{minioadmin} can create buckets,
delete buckets, and read every object in every bucket. The service
needs \code{PutObject}, \code{GetObject} and \code{HeadBucket} on one
bucket. Least privilege here is a MinIO user with a policy scoped to
\code{arn:aws:s3:::ts-submissions/*}; on AWS it is an IAM role attached
to the pod's ServiceAccount (IRSA) so that no static key exists at all.
An interviewer who asks ``how does your service authenticate to S3?''
is listening for whether you say ``an IAM role'' or ``a key in an env
var'' --- and, if the latter, whether you know it is the weaker answer.

\section{What breaks}

\begin{pitfall}[the service that will not start without its bucket]
Symptom: course-service in \code{CrashLoopBackOff} on a fresh cluster;
logs end in \code{Error creating bean with name 's3StorageService'} with
a \code{Connection refused} several causes down. Cause: the
\code{@PostConstruct} bucket check treats any failure other than
``bucket missing'' as fatal, and MinIO's pod is still pulling its image.
The kubelet's restart back-off eventually lands on a moment when MinIO
is ready, so it self-heals --- slowly. Fix: catch
\code{SdkClientException} too, log a warning, and let the readiness
probe (Chapter~17) keep traffic away until the first successful
\code{headBucket}; a storage dependency should degrade the file feature,
not the service.
\end{pitfall}

\begin{pitfall}[orphans]
Symptom: the bucket holds more objects than the \code{submission} table
has non-null keys; nobody notices for a year; the 5\,GB claim fills.
Cause: the dual write in \code{submit}. Upload first, insert second: a
failed insert orphans an object. Insert first, upload second: a failed
upload leaves a row pointing at nothing, and the student sees ``file
attached'' with a dead link. Neither order is safe on its own. The
standard resolution is to make one side idempotent and the other side
reconcilable: insert the row with status \code{PENDING\_FILE} and the
chosen key, upload, then flip the status to \code{SUBMITTED} in a second
transaction; a nightly job deletes objects under keys whose row is still
\code{PENDING\_FILE} after an hour, and lists the prefix for keys no row
names at all. Section~\ref{sec:presigned-put} folds this into the new
flow, where it comes almost free.
\end{pitfall}

\section{What breaks at 3\,a.m.}

MinIO's pod is gone --- node reboot, claim not yet reattached. Walk the
endpoints:

\begin{itemize}
  \item \textbf{Text and link submissions} still work: \code{upload} is
  only called when a file is present, and \code{save} touches only
  Postgres.
  \item \textbf{File submissions} fail with a 500 after the SDK's retry
  budget --- several seconds of a held thread each. Under load these
  threads accumulate and starve the text submissions that would have
  worked. This is the Chapter~17 argument for a bulkhead: a separate,
  bounded executor for storage calls.
  \item \textbf{Listing submissions} works and even returns download
  links. Presigning is pure computation; no request is made. The links
  fail only when clicked, which is the correct failure --- the metadata
  is true, the bytes are temporarily unavailable.
  \item \textbf{A course-service restart} during the outage does not
  come back up, for the reason in the pitfall above. An outage of the
  file store has become an outage of grading, enrollment, and everything
  else the service does.
\end{itemize}

Now the other way: Postgres is down and MinIO is fine. Everything fails,
including uploads that would have succeeded at the storage layer ---
and each one that reached \co{1} in \code{submit} before failing at
\co{2} left an orphan behind. A 3\,a.m. database blip produces a small
pile of unowned objects, which is the concrete reason the reconciliation
job exists.

\section{Enhancing the resolution: presigned PUT}
\label{sec:presigned-put}

The upgrade moves the bytes off the service and closes the orphan
window at the same time. Three calls replace one:

\begin{enumerate}
  \item \code{POST /assignments/\{id\}/submissions/upload-url} --- the
  service checks the student may submit, chooses the key, signs a PUT,
  and inserts a \code{PENDING\_FILE} row.
  \item The browser PUTs the file to MinIO with that URL.
  \item \code{POST /assignments/\{id\}/submissions} with the key --- the
  service confirms the object exists with a \code{HEAD}, then completes
  the row.
\end{enumerate}

The storage service gains two methods. Both are complete:

\begin{lstlisting}[language=Java]
/** Sign a PUT for a key we chose; the browser must send exactly this
 *  content type, because it is part of the signature. */
public String presignedUploadUrl(String key, String contentType,
                                 long contentLength) {
    PutObjectRequest put = PutObjectRequest.builder()
            .bucket(bucket)
            .key(key)
            .contentType(contentType)
            .contentLength(contentLength)                  (*@\co{1}@*)
            .build();
    PutObjectPresignRequest req = PutObjectPresignRequest.builder()
            .signatureDuration(Duration.ofMinutes(10))
            .putObjectRequest(put)
            .build();
    return s3Presigner.presignPutObject(req).url().toString();
}

/** Confirm an object exists and report its size; empty if missing. */
public Optional<Long> objectSize(String key) {
    try {
        HeadObjectResponse head = s3Client.headObject(
                HeadObjectRequest.builder().bucket(bucket).key(key).build());
        return Optional.of(head.contentLength());
    } catch (NoSuchKeyException e) {                        (*@\co{2}@*)
        return Optional.empty();
    }
}
\end{lstlisting}

\begin{description}
  \item[\co{1}] Signing \code{Content-Length} and \code{Content-Type}
  turns them into a contract: a browser that sends a different size or
  type gets \code{403}. That is how the 20\,MB limit survives the move
  out of the service.
  \item[\co{2}] \code{HEAD} costs one round trip and no bytes. It is the
  cheapest possible proof that step~2 happened.
\end{description}

The submission service splits its old method in two. The key prefix is
built from server-side facts, so a client cannot claim someone else's
upload:

\begin{lstlisting}[language=Java]
public UploadTicket prepareUpload(Long assignmentId, Long studentUserId,
                                  String fileName, String contentType,
                                  long size) {
    Assignment assignment = requireOpenAssignment(assignmentId);
    Student student = requireEnrolled(assignment, studentUserId);
    if (size > MAX_FILE_BYTES) {
        throw new ApiException(413, "File exceeds 20 MB");
    }
    String key = "submissions/%d/%d/%s-%s".formatted(
            assignmentId, student.getId(), UUID.randomUUID(),
            safeName(fileName));                            (*@\co{1}@*)
    Submission pending = Submission.builder()
            .assignment(assignment).student(student)
            .attemptNumber(nextAttempt(assignmentId, student.getId()))
            .fileObjectKey(key).fileName(fileName)
            .status("PENDING_FILE")                         (*@\co{2}@*)
            .build();
    submissionRepository.save(pending);
    String url = s3StorageService.presignedUploadUrl(key, contentType,
            size);
    return new UploadTicket(pending.getId(), key, url);
}

public SubmissionResponse completeUpload(Long submissionId,
                                         Long studentUserId) {
    Submission s = findOrThrow(submissionId);
    if (!s.getStudent().getUserId().equals(studentUserId)
            || !"PENDING_FILE".equals(s.getStatus())) {
        throw new ApiException(409, "Nothing to complete");
    }
    long size = s3StorageService.objectSize(s.getFileObjectKey())
            .orElseThrow(() -> new ApiException(400,
                    "Upload not found; request a new upload URL")); (*@\co{3}@*)
    if (size > MAX_FILE_BYTES) {
        throw new ApiException(413, "File exceeds 20 MB");
    }
    s.setStatus("SUBMITTED");
    return toResponse(submissionRepository.save(s));
}
\end{lstlisting}

\begin{description}
  \item[\co{1}] \code{safeName} strips path separators and
  non-ASCII; the display name stays in \code{file\_name}.
  \item[\co{2}] The row exists \emph{before} the bytes. A crash anywhere
  after this point leaves a \code{PENDING\_FILE} row that the cleanup
  job can find, instead of an object nobody can find.
  \item[\co{3}] The only way to reach \code{SUBMITTED} is through a
  successful \code{HEAD}. The row never claims a file that is not there.
\end{description}

The cleanup job is now a query, not a bucket scan:

\begin{lstlisting}[language=Java]
@Scheduled(cron = "0 15 * * * *")   // hourly, at :15
public void reapAbandonedUploads() {
    List<Submission> stale = submissionRepository
            .findByStatusAndSubmittedAtBefore("PENDING_FILE",
                    LocalDateTime.now().minusHours(1));
    for (Submission s : stale) {
        s3StorageService.deleteQuietly(s.getFileObjectKey()); // idempotent
        submissionRepository.delete(s);
    }
}
\end{lstlisting}

MinIO must allow the browser's cross-origin PUT. The bucket's CORS rule
is set once with the MinIO client:

\begin{lstlisting}[language=bash]
$ mc cors set local/ts-submissions cors.json
# cors.json: allow PUT from the frontend origin, expose ETag
{ "CORSRules": [ { "AllowedOrigins": ["http://localhost:3000"],
    "AllowedMethods": ["PUT", "GET"],
    "AllowedHeaders": ["Content-Type", "Content-Length"],
    "ExposeHeaders": ["ETag"], "MaxAgeSeconds": 3600 } ] }
\end{lstlisting}

And the frontend side is a plain \code{fetch}; no axios instance, no
gateway, no JWT --- the URL \emph{is} the credential:

\begin{lstlisting}
const ticket = await api.post(
  `/assignments/${id}/submissions/upload-url`,
  { fileName: file.name, contentType: file.type, size: file.size });

const put = await fetch(ticket.url, {
  method: 'PUT',
  headers: { 'Content-Type': file.type },   // must match the signature
  body: file,
});
if (!put.ok) throw new Error(`upload failed: ${put.status}`);

await api.post(`/submissions/${ticket.submissionId}/complete`);
\end{lstlisting}

\begin{decision}[presigned PUT versus proxied upload]
Proxy when the service must see the bytes before they are stored ---
virus scanning, image re-encoding, content-type sniffing --- or when the
storage endpoint cannot be reached from the client at all. Presign when
files are large or many and the service adds nothing to the transfer:
the thread pool, the ingress body limit, and the double network hop all
disappear. The project's case is the second, with one caveat that
decides the order of work: presigned uploads require a browser-reachable
MinIO origin, which the cluster does not have yet. The pitfall above and
this enhancement share a prerequisite --- a public files hostname ---
and should be scheduled together.
\end{decision}

\section{Outside the project}

On AWS the same code runs against S3 with the endpoint override removed
and static keys replaced by an IAM role; CloudFront signed URLs play the
presigner's role with a CDN in front, and \emph{S3 POST policies} are the
older, browser-form-friendly cousin of presigned PUT, with a
\code{content-length-range} condition built in. Google Cloud Storage and
Azure Blob Storage have the same three ideas under other names ---
signed URLs, SAS tokens, resumable uploads. Production MinIO runs as a
distributed set with erasure coding across drives and nodes; a single
container with one claim is a development shape, and the SDK cannot tell
the difference, which is the point of the standard. The interview
version of this chapter is one question --- ``a user uploads a 2\,GB
video; walk me through it'' --- and the senior answer names presigned or
multipart upload, a size and type contract in the signature, a
confirmation step, an orphan reaper, and a lifecycle rule for
incomplete parts, in about that order.

\section{Questions from real interviews}

\subsection*{Q: Users started uploading 1\,GB videos and the service runs out of memory. Fix it.}

First establish that the bytes are the problem, not the metadata: a heap
dump during an upload shows multipart buffers or a
\code{ByteArrayOutputStream}, and the pod's memory limit
(Chapter~17) is what kills it. The short-term fix is to make sure Spring
spools to disk (\code{file-size-threshold} small, a temp volume with room)
and stream to storage with a known \code{Content-Length}, as
\code{S3StorageService.upload} already does --- that gets you to
hundreds of MB without heap growth, but still one thread and one
temp file per upload. The real fix is to take the service out of the
path: presigned PUT for anything under about 100\,MB, and multipart
upload with \code{S3TransferManager} or browser-side parts above it,
with a lifecycle rule aborting incomplete uploads after seven days. Then
the service's memory is independent of file size by construction, and
\code{max-file-size} becomes a signed \code{Content-Length} instead of
a servlet setting.

\subsection*{Q: We need to serve ten thousand downloads per second. How, without touching the service?}

Presigned GET already means the service does no I/O per download, only
a signature --- microseconds, no network. Ten thousand per second is
therefore a storage and edge question. Object stores scale reads
horizontally by key prefix distribution; a CDN in front (CloudFront with
signed URLs, or any CDN with origin-pull to the bucket) serves repeated
objects from cache and absorbs the fan-out. Keep the signing cheap by
caching nothing --- signatures are stateless --- and set the expiry to
match the CDN's cache window so a cached object does not outlive its
authorization by much. The trap is signing on every list call:
\code{toResponse} in the project presigns for every row of every
listing, which is fine at 15 minutes' expiry but means a listing of a
thousand submissions computes a thousand HMAC chains. Sign lazily, on a
dedicated ``get link'' endpoint, once volume justifies it.

\subsection*{Q: An object exists in the bucket but the row is missing, or the row exists and the object does not. How do you reconcile?}

Say first that both are expected outcomes of a dual write and that the
fix is design, not a script. Design: write the row in a
\code{PENDING} state with the key before the upload, confirm with
\code{HEAD} before marking it complete, and never let a row claim a
file it has not verified. Reconciliation then has two cheap halves. Rows
stuck in \code{PENDING} older than the presign expiry are deleted along
with their objects --- a query, not a scan. Objects with no row are found
by listing the prefix and anti-joining against the table, which you can
afford nightly because the prefix is per assignment; on AWS an S3
Inventory report replaces the listing. Whatever you do, make the object
delete idempotent so the job can be rerun.

\subsection*{Q: How do you secure a presigned URL?}

Four dials, and a senior answer names all of them. \emph{Scope}: sign
one method on one key; a GET URL cannot be used to PUT, and a PUT URL
cannot be used on a different key because the canonical URI is signed.
\emph{Expiry}: minutes, not days; the project uses 15 for downloads and
10 would do for uploads --- long enough for a slow connection, short
enough that a leaked link in a chat log is dead by the time it is
noticed. \emph{Content contract}: sign \code{Content-Type} and
\code{Content-Length} on uploads so the client cannot change what it
declared. \emph{Identity of the signer}: the credentials behind the
signature should be a least-privilege user on that one bucket, so even a
compromised signer cannot list or delete. And say the thing people
forget: a presigned URL is a bearer credential --- log it nowhere,
return it over HTTPS only, and never put it in a redirect that a proxy
might cache.

\subsection*{Q: Migrate from MinIO to AWS S3 with zero downtime.}

Because the code speaks the S3 API, the migration is a data move plus a
config change, and the trick is ordering. Set up bucket replication or
run \code{mc mirror --watch} from MinIO to the S3 bucket until the lag
is seconds. Switch \emph{writes} first by changing
\code{APP\_S3\_ENDPOINT} and credentials (an IAM role, now) on a rolling
restart --- new objects land in S3, mirroring keeps copying anything
still landing in MinIO from pods not yet restarted. Switch the
\emph{public endpoint} last, after a final mirror pass, so download
links are only ever signed for a host that has the object. Keys do not
change, so the database needs no migration. Keep MinIO read-only for a
week as a fallback, then decommission. The thing that breaks is
addressing: virtual-host style becomes available, path-style keeps
working, and any code that concatenated URLs by hand instead of asking
the SDK is where you will find the bug.

\subsection*{Q: Storage cost is growing every month. What do you do?}

Measure by prefix first --- which assignments, which years. Then use the
store's own tools rather than application logic: a lifecycle rule that
transitions objects older than a term to an infrequent-access tier, and
older than the retention policy to archive or deletion; an abort rule
for incomplete multipart uploads, which are invisible and billed; and
versioning off, or with expiry, unless a compliance need says otherwise.
In the application, stop storing what you can regenerate (thumbnails,
transcodes) and deduplicate by content hash if students upload the same
template. On this project the numbers are small --- 250 uploads of
20\,MB fill the 5\,GB claim --- but the same rule set is what keeps an
S3 bill flat at a thousand times the size.

\begin{quiz}
\begin{enumerate}
  \item Two beans, \code{S3Client} and \code{S3Presigner}, share
  credentials and region and differ in one line. Name the line, and
  explain what goes wrong if the presigner used the internal endpoint.
  \item A colleague suggests keeping one endpoint and having the gateway
  rewrite \code{minio:9000} to the public host in the JSON response.
  Using the structure of a Version 4 signature, explain why the rewritten
  link returns 403.
  \item In \code{SubmissionService.submit}, the upload happens inside a
  \code{@Transactional} method. State what the transaction does and does
  not cover, and describe the artifact left behind by a failed
  \code{save}.
  \item The key layout is \code{submissions/\{assignment\}/\{student\}/\{uuid\}-\{name\}}.
  Give one operation the layout makes cheap and one it makes expensive,
  and say why the UUID makes bucket versioning unnecessary for safety.
  \item MinIO goes down at 3\,a.m. Which of these still succeed:
  submitting a text answer, listing submissions, opening a download
  link, restarting course-service? Explain each in one clause.
  \item In the presigned-PUT flow, an attacker obtains a valid upload
  URL from their own session and tries to attach that object to another
  student's submission. Identify the two checks that stop them.
\end{enumerate}
\end{quiz}
